\section{Plano de Validação Matemática}

\subsection{Propriedades Estruturais do Modelo}

As propriedades estruturais do modelo são fundamentais para garantir a consistência física e matemática do sistema. A matriz de inércia $ \mathbf{M}(\vec{q}) $ deve ser simétrica e definida-positiva para todo $ \vec{q} $. Além disso, o modelo deve obedecer à propriedade de passividade, ou seja, para todo vetor $ \vec{y} $, a seguinte identidade deve ser satisfeita:

\begin{equation}
\vec{y}^T \left( \dot{\mathbf{M}}(\vec{q}) - 2\mathbf{C}(\vec{q}, \dot{\vec{q}}) \right) \vec{y} = 0
\label{eq:passividade}
\end{equation}

Essa é uma característica clássica de modelos Lagrangianos. Também se espera a conservação de energia na ausência de entradas e dissipação, o que leva à seguinte equação:

\begin{equation}
\dot{\mathcal{E}} = \frac{d}{dt}(T + V) = 0
\label{eq:conservacao_energia}
\end{equation}

\subsection{Consistência de Jacobianos}

Para garantir a validade da cinemática diferencial implementada.
Neste trabalho, a consistência dos Jacobianos foi avaliada por diferenças finitas.
A Jacobiana $ \mathbf{J} $ pode ser aproximada por:

\begin{equation}
\mathbf{J} \approx \frac{\vec{x}_E(\vec{q} + \Delta \vec{q}) - \vec{x}_E(\vec{q})}{\Delta \vec{q}}, \qquad \left\| \mathbf{J} - \mathbf{J}_{\mathrm{FD}} \right\| \leq \varepsilon_J
\label{eq:jacobiano_finitas}
\end{equation}

Esse procedimento deve ser repetido para cada elo e articulação utilizada no modelo.

\subsection{Checagens de Unidades e Escalas}

A coerência dimensional do modelo é verificada com as seguintes expressões:

\begin{itemize}
\item $ [\mathbf{M}] = \mathrm{kg} \cdot \mathrm{m}^2 $
\item $ [\mathbf{C}] = \mathrm{kg} \cdot \mathrm{m}^2 / \mathrm{s} $
\item $ [\vec{g}] = \mathrm{N} \cdot \mathrm{m} $ (ou $ \mathrm{N} $ em coordenadas prismáticas)
\end{itemize}

Neste trabalho, quando houve problemas de condicionamento, foi aplicada a normalização para reescalar variáveis.

\subsection{Vínculos e Estabilização}

Durante a simulação, o erro de vínculo $ \| \Phi(\vec{q}) \| $ e o resíduo de Baumgarte devem permanecer abaixo dos limites $ \varepsilon_\Phi $ e $ \varepsilon_B $. No caso de uso de quaternion, é necessário monitorar a normalização:

\begin{equation}
\left\| \vec{q}_b \right\| - 1 \leq \varepsilon_q
\label{eq:erro_quaternion}
\end{equation}

% \subsection{Linearização e Discretização}

% A validade do modelo linearizado foi avaliada por meio da análise do espectro da matriz $A$ e da controlabilidade local em torno do ponto de operação. Além disso, comparou-se a resposta do modelo linear discreto com o não linear sob pequenas excitações, utilizando a formulação

% \begin{equation}
% \vec{x}_{k+1} = \mathbf{A}_d \vec{x}_k + \mathbf{B}_d \vec{u}_k, \quad
% \mathbf{A}_d = e^{\mathbf{A}T_s}, \quad
% \mathbf{B}_d = \int_0^{T_s} e^{\mathbf{A}\tau} d\tau \, \mathbf{B}
% \label{eq:discretizacao_ZOH}
% \end{equation}


\subsection{Energia e Potência}

Na ausência de dissipação e entradas externas, espera-se que a energia total do sistema permaneça aproximadamente constante. Com a inclusão de dissipação e entradas generalizadas, o balanço de potência segue a relação:
\begin{equation}
\mathcal{E}(t) \approx \mathcal{E}(0), \quad \text{com variação relativa} \leq \varepsilon_E
\label{eq:energia_constante}
\end{equation}

Com dissipação $ \mathbf{D} $ e entradas $ \vec{\tau} $, o balanço de energia esperado é:

\begin{equation}
\dot{\mathcal{E}} \approx \vec{\tau}^T \dot{\vec{q}} - \dot{\vec{q}}^T \mathbf{D} \dot{\vec{q}}
\label{eq:balanco_energia}
\end{equation}

\subsection{Sensibilidade Paramétrica}

A robustez do sistema foi investigada mediante variações plausíveis nas massas, centros de massa e tensores de inércia. Avaliaram-se os impactos nas saídas, como desvios de posição do efetuador, e no esforço de controle requerido.

\subsection{Tolerâncias Sugeridas}

Para fins de validação numérica, adotaram-se tolerâncias típicas da ordem de:

\begin{itemize}
\item $ \varepsilon_J \sim 10^{-3} $;
\item $ \varepsilon_\Phi \sim 10^{-6} $;
\item $ \varepsilon_E \sim 10^{-3} $;
\item $ \varepsilon_q \sim 10^{-6} $.
\end{itemize}

\subsection{Observações para a Implementação}

% Verifica-se a convergência do solver integrando com diferentes passos e tolerâncias. O condicionamento numérico das matrizes $ \mathbf{M} $ e $ \mathbf{J} $ também deve ser monitorado. Por fim, recomenda-se comparar a cinemática e dinâmica com uma segunda fonte confiável, como o modelo do \texttt{rigidBodyTree} em posturas de teste.
A convergência numérica foi avaliada mediante integração com diferentes passos e tolerâncias, monitorando o condicionamento das matrizes $M$ e $J$. A consistência da cinemática e da dinâmica também foi verificada por comparação com modelos de referência, como o rigidBodyTree em posturas de teste
